\documentclass[11pt,letterpaper]{article}
\usepackage[utf8]{inputenc} %Codificacion del texto (ISO Latin1 encoding)

\usepackage{fancyhdr} %Permite acomodar a tu gusto la parte de arriba y
\usepackage[spanish]{babel} %Permite definir el idioma del dcumento
\usepackage{graphicx} %Permite exportar imagenes en formato eps
\usepackage{url} %Tipo de fuente para correos y paginas
\usepackage{hyperref}
\usepackage{pgf}
\usepackage{fleqn}
\usepackage{amssymb}
\usepackage{fancyvrb}
\usepackage{sectsty}
\usepackage{makeidx}
\usepackage{colortbl} %Permite colocar colores a las tablas
\usepackage{booktabs}
%Margenes%
\parskip 1mm %Espacio entre parrafos

\setlength{\topmargin}{0pt}

\oddsidemargin	0.5cm  % Ancho Letter 21,59cm
\evensidemargin 0.5cm  % Alto  Letter 27,81cm
\textwidth	15.5cm
\textheight	21.0cm
\headsep	4 mm
\parindent	0.5cm
\pagestyle{fancyplain}
\lhead{Taller de Gestión de Proyectos Informáticos} %Parte superior izquierda
\rhead{\bf \it Sintesis} %Parte superior derecha
\lfoot{\it Principios y Conceptos de la Gestión de Proyectos }
\cfoot{} %Parte inferior central
\rfoot{\bf \thepage} %Parte inferior derecha
\renewcommand{\footrulewidth}{0.4pt} %Linea de separacion inferior

\newtheorem{theorem}{Theorem}
\newtheorem{acknowledgement}[theorem]{Acknowledgement}
\newtheorem{algorithm}[theorem]{Algorithm}
\newtheorem{axiom}[theorem]{Axiom}
\newtheorem{case}[theorem]{Case}
\newtheorem{claim}[theorem]{Claim}
\newtheorem{conclusion}[theorem]{Conclusion}
\newtheorem{condition}[theorem]{Condition}
\newtheorem{conjecture}[theorem]{Conjecture}
\newtheorem{corollary}[theorem]{Corollary}
\newtheorem{criterion}[theorem]{Criterion}
\newtheorem{definition}[theorem]{Definition}
\newtheorem{example}[theorem]{Example}
\newtheorem{exercise}[theorem]{Exercise}
\newtheorem{lemma}[theorem]{Lemma}
\newtheorem{notation}[theorem]{Notation}
\newtheorem{problem}[theorem]{Problem}
\newtheorem{proposition}[theorem]{Proposition}
\newtheorem{remark}[theorem]{Remark}
\newtheorem{solution}[theorem]{Solution}
\newtheorem{summary}[theorem]{Summary}
\newenvironment{proof}[1][Proof]{\noindent\textbf{#1.} }{\ \rule{0.5em}{0.5em}}

\newcommand{\primaria}[1]{
	\textbf{\underline{#1}}
}

\newcommand{\foranea}[1]{
	\textbf{\textsl{#1}}
}

\newcommand{\primyfor}[1]{
	\underline{\foranea{#1}}
}

\makeatletter
\newcommand\subsubsubsection{\@startsection {paragraph}{1}{\z@}%
                                   {-3.5ex \@plus -1ex \@minus -.2ex}%
                                   {1.5ex \@plus.2ex}%
                                   {\normalfont\bfseries}}
\newcommand\subsubsubsubsection{\@startsection {subparagraph}{1}{\z@}%
                                   {-3.5ex \@plus -1ex \@minus -.2ex}%
                                   {1.5ex \@plus.2ex}%
                                   {\normalfont\bfseries}}


\makeatother
%\makeindex
\bibliographystyle{plain}

\begin{document}

%%%%%%%%%%%%%%%%%%%%%%%%%%
%Definicion de la portada%
%%%%%%%%%%%%%%%%%%%%%%%%%%
\begin{titlepage}
    \begin{center}
	%\begin{tabular}{ccc}
	\begin{tabular}{c}
		\includegraphics[width=0.9\textwidth]{img/logos}
		
	   % \includegraphics[width=3cm]{img/utfsm}
	   % & 
	   % \hspace{-0.2cm}
	   % \begin{tabular}{c}
	   % Universidad Técnica Federico Santa María \\ \hline
	   % \vspace{0.2cm}
	   % Departamento de Informática\\
	   % \vspace{1.2cm}
	   % \end{tabular}
	   % \hspace{0.2cm}
	   % &
        %    \includegraphics[width=2cm]{img/di}
	\end{tabular}

	\vspace{4cm}
	%Titulo del Documento
	\begin{tabular}{c}
		\vspace{3cm}
		\Large{\sc{Seminario de Modelos y Métodos Cuantitativos}}\\
		\huge{\sc{Tarea 2}}\\\\
		%\includegraphics[scale=0.7]{img/portada} \\\\
	\end{tabular}

    \vspace{5cm}
	\begin{tabular}{lr}
			\textbf{Alumnos} & \\
							 & \\
         	\normalsize{Cristián Maureira Fredes} & \url{cmaureir@csrg.inf.utfsm.cl}\\
         	\normalsize{Gabriel Zamora Nelson} & \url{gzamora@csrg.inf.utfsm.cl}\\

							 & \\
			\textbf{Profesor} & \\
							 & \\
         	\normalsize{Andrés Moreira} & \url{amoreira@inf.utfsm.cl}\\
	\end{tabular}

\vspace{2cm}

	%Fecha
    \normalsize{\sc{\today}}\\
    %\normalsize \textbf{Fecha de Entrega:} & {14 de Noviembre del 2010}\\
    \end{center}
\end{titlepage}

\tableofcontents
\newpage

%\section{Introducción}
%\label{sec:introduccion}
%\frame
{
\frametitle{Introducción}
\begin{itemize}
	\item Motivación del problema.
	\item Sobre la técnica.
	\item Implementación inicial.
	\item Sintonización.
	\item \blue{Control}.
	\item Resultados.
\end{itemize}
}



\section{Administración del plan de riesgo}

Es el proceso de definir como realizar las actividades de gestión de riesgo del proyecto.
En el se debe realizar una planificación cuidadosa y bien definida, ya que
ésta mejora las probabilidades de triunfar en los 5 otros procesos de la
administración del riesgo. Por ello, es que éste proceso comienza cuando el proyecto es concebido.

La administración del plan de riesgo requiere considerar los diferentes planes
y factores que posee la empresa. Por ello, hay que considerar los alcances del
proyecto, los planes de administración de costos, tiempos y comunicaciones,
los factores del ambiente empresarial y la información organizacional de los
procesos, para poder definir bien las categorías de riesgo, definiciones en
común de términos y conceptos, formatos de presentación de riesgos, templates
estandarizados, roles y responsabilidades, niveles de autoridad, etc.

Para poder llevar a cabo de buena forma el proceso, es recomendable tener en
cuenta diferentes técnicas y herramientas, dentro de las que se encuentra el
realizar reuniones de análisis y elaboración del plan de administración de
riesgo. En estas reuniones se requieren todos los involucrados del proyecto, y
se busca definir las actividades que se seguirán para controlar los posibles
riesgos y sus costos respectivos, las medidas de contingencia y asignación de
responsabilidades. Éstas se deben categorizar por nivel, probabilidad,
impacto, etc. del riesgo, y generar los templates necesarios para el correcto
control de los mismos.

Al final del proceso se obtendrían los documentos que definen la metodología,
roles y responsabilidades, presupuesto, periodos de evaluación, categorías de
riesgos, definiciones de probabilidad y impacto de los riesgos, la matriz de
probabilidad e impacto, la tolerancia de los stakeholders, los formatos de
reporte y por último la documentación y seguimiento del proceso.
\subsection{Valores de entrada}

\begin{itemize}
    \item Alcances del proyecto.
    \item Plan de la administración de costos.
    \item Plan de administración de tiempos.
    \item Plan de administración de Comunicaciones.
    \item Factores del ambiente empresarial.
    \item Información organizacional de los procesos.
    \begin{itemize}
        \item Categorías de riesgo.
        \item Definiciones en común de términos y conceptos.
        \item Formatos de presentación de riesgos.
        \item Templates estandarizados.
        \item Roles y responsabilidades.
        \item Niveles de autoridad para la toma de decisiones.
        \item Lecciones aprendidas.
        \item Registros de los stakeholders.
    \end{itemize}
\end{itemize}

\subsection{Técnicas y Herramientas}
\begin{itemize}
    \item Reuniones de análisis y elaboración del plan de administración de riesgo.
    \begin{itemize}
        \item Requieren de todos los involucrados del proyecto.
        \item Se busca definir las actividades para la administración del
        riesgo.
        \item Se calculan los costos respectivos.
        \item Revisiones de las medidas de contingencia.
        \item Asignación de responsabilidades.
        \item Categorización de riesgos por nivel, probabilidad, impacto, etc.
        \item Generación de templates.
    \end{itemize}
\end{itemize}

\subsection{Valores de Salida}
\begin{itemize}
    \item Metodología.
    \item Roles y responsabilidades.
    \item Presupuesto.
    \item Periodos de evaluación.
    \item Categorías de riesgos.
    \item Definiciones de probabilidad y impacto de los riesgos.
    \item Matriz de probabilidad e Impacto.
    \item Tolerancia de los stakeholders.
    \item Formatos de reporte.
    \item Documentación y seguimiento.
\end{itemize}

\section{Identificar Riesgos}
Es el proceso de identificar los riesgo que puedan afectar al proyecto,
documentando sus características principales. Esto es iterativo, ya que nuevos
riesgos pueden evolucionar o aparecer en el transcurso del proyecto. Por ello,
debe involucrar a los miembros del proyecto para inculcar un sentido de
responsabilidad en ellos.

En este proceso, se deben considerar los planes realizados anteriormente,
junto con los costos y duración de cada una de las actividades propuestas y el
lineamiento base sobre los alcances del proyecto. Por ello, los registros con
información sobre los stakeholders y los planes de administración de costos,
tiempos y calidad, junto con los documentos del proyecto, deben ser revisados
exhaustivamente para poder identificara los principales riesgos.

En esta etapa, se debería revisar la documentación que se tenga del proyecto,
y utilizar técnicas de recopilación de la información como Brainstorming,
técnica de Delphi, entrevistas, etc. para poder contar con una buena
información de todos los niveles de trabajo del proyecto. El realizar análisis
con listas de chequeo, analizar las suposiciones y el desarrollo de diagramas
de causa y efecto, flujo del sistema y influencia pueden ayudar a la
identificación de riesgos. Por último, vale la pena nombrar el análisis SWOT
(FODA) y el juicio de expertos para una revisión más profunda de los riesgos.

El resultado de la identificación de riesgos es una documentación con el
registro de riesgos, identificados apropiadamente y con acciones potenciales
para la mitigación de cada uno de ellos.
\subsection{Valores de Entrada}

\begin{enumerate}
    \item Plan de administración de Riesgos.
    \item Costos estimados de las actividades.
    \item Duración estimada de las actividades.
    \item Lineamiento base sobre los alcances del proyecto.
    \item Registros de los stakeholders.
    \item Plan de administración de Costos.
    \item Plan de administración de Tiempos.
    \item Plan de administración de Calidad.
    \item Documentos del proyecto
    \begin{itemize}
        \item Diagramas de red, lineamientos base, eficiencia de trabajo, etc.
    \end{itemize}
    \item Factores del ambiente empresarial.
    \begin{itemize}
        \item Información pública, bases de datos, estudios académicos,
        benchmarking, estudios de la industria, etc.
   
    \end{itemize}
    \item Información organizacional de los procesos.
    \begin{itemize}
        \item Información del proyecto, procesos de control del proyecto,
        templates, lecciones aprendidas, etc.
    \end{itemize}

\end{enumerate}


\subsection{Técnicas y Herramientas}

\begin{enumerate}
    \item Revisiones de la documentación.
    \item Técnicas para la recopilación de la información.
    \begin{itemize}
        \item Brainstorming.
        \item Técnica de Delphi.
        \item Entrevistas.
        \item Análisis de las fuentes del problema.
    \end{itemize}
    \item Análisis con listas de chequeo.
    \item Análisis de las suposiciones.
    \item Técnicas de desarrollo de diagramas.
    \begin{itemize}
        \item Diagramas de causa y efecto, flujo del sistema y influencia.
    \end{itemize}
    \item Análisis SWOT (FODA).
    \item Juicio de Expertos.
\end{enumerate}

\subsection{Valores de Salida}
\begin{enumerate}
    \item Registro de riesgos
    \begin{itemize}
        \item Lista de riesgos identificados.
        \item Lista de acciones potenciales.
    \end{itemize}
\end{enumerate}

\section{Análisis cualitativas del riesgo}
Realizar "análisis de riesgos cualitativo" es el proceso de
priorización de riesgos para realizar análisis adicionales o
determinar el grado de evolución de los actos de una organización,
asociados a una cierta probabilidad de ocurrencia y su impacto.

Las organizaciones pueden mejorar el rendimiento de los proyectos
concentrándose en riesgos de mayor prioridad.
No hay que dejar de lado que existe una Probabilidad de ocurrencia
para cada riesgo, tanto de mayor o menos prioridad.

La presente evaluación refleja también la actitud del equipo del proyecto y
los stakeholders a los riesgos estudiados.
Los tiempos críticos de las acciones relacionadas al riesgo pueden
magnificar la importancia de un riesgo.

La evaluación de la calidad de la información de un proyecto
puede ayudar a la claridad de la evaluación de riesgos.
Realizar esto es una forma rápida y de bajo costo para establecer
las prioridades para las respuestas de un plan de riesgo.

Se realiza durante todo el ciclo de vida del proyecto,
ya que hay ciertos cambios a lo largo del desarrollo del proyecto
que pueden llevar a poseer nuevos riesgos o a eliminar algunos.

\subsection{Entradas}
\begin{enumerate}
	\item Registro de Riesgos.
	\item Plan de gestión de riesgos. (roles, responsabilidades, planificación
		de actividades, categorías de riesgos, definiciones, definición de probabilidad
		e impacto, matriz de impacto, tolerancia de riesgo de los stakeholders.
	\item Declaración del ámbito del proyecto. (proyectos parecidos, tienen los mismos
		riesgos y los que son nuevos tienen más incertidumbre)
	\item Activos de procesos de organización. (Información de proyectos similares
		anteriores, estudios por especialistas, bases de datos)
\end{enumerate}

\subsection{Herramientas y Técnicas}
\begin{enumerate}
	\item Probabilidad de riesgo y evaluación del impacto.
	\item Probabilidad y matriz de impacto. (oportunidades y amenazas, ranking)
	\item Datos de los riesgo en la evaluación de calidad.
	\item Categorización de riesgos.
	\item Evaluación de urgencia de riesgo.
	\item Juicio Experto.
\end{enumerate}


\subsection{Salidas}
\begin{enumerate}
	\item Actualizaciones de registro de riesgos
        Se realiza un registro cuando se identifican los procesos de riesgos.
        Esta información se actualiza con la información del Análisis de Riesgo Cualitativo.
        Cada actualización contiene:
		\begin{itemize}
            \item Ranking relativo o lista de prioridades de los riesgos del proyecto.
			\item Riesgos agrupados por categorías.
			\item Causas de riesgos o áreas de proyecto que requieren una atención particular.
			\item Lista de riesgos que requieren respuesta cercana al equipo de trabajo.
			\item Lista de riesgos para un análisis y respuesta adicional.
			\item Lista de seguimiento de riesgos de baja prioridad.
			\item Tendencias en el resultado del análisis  de riesgos cualitativo.
		\end{itemize}
\end{enumerate}


\section{Análisis cuantitativos del riesgo}
\frame
{
\frametitle{Análisis Cuantitativo de Riesgos}
\begin{itemize}
	\item Este proceso usa \emph{técnicas y métodos} \textbf{estadísticos} para:
	\begin{itemize}
		\item Cuantificar los \emph{posibles resultados} del proyecto y sus probabilidades
		\item Evaluar la \emph{probabilidad de lograr} los \textbf{objetivos} específicos del proyecto.
		\item Identificar los riesgos que requieren una \textbf{mayor atención}.
	\end{itemize}
\end{itemize}
\begin{columns}
	\begin{column}{0.6\textwidth}
	\end{column}
	\begin{column}{0.4\textwidth}
		\includegraphics[width=3cm]{img/quantity}
	\end{column}
\end{columns}
}

\frame
{
\frametitle{Análisis Cuantitativo de Riesgos}
\begin{block}{Definición}
\begin{itemize}
	\item Analiza el efecto de esos riesgos y les asigna una \textbf{calificación numérica}.
	\item También presenta un método cuantitativo para \textbf{tomar decisiones} en caso de \emph{incertidumbre}. 
\end{itemize}
\end{block}
}


\frame
{
\frametitle{Análisis Cuantitativo de Riesgos}
\framesubtitle{Entradas}
\begin{columns}
	\begin{column}{0.8\textwidth}
		\begin{itemize}
			\item<1-> Registro de Riesgos.
			\item<2-> Plan de administración de \textbf{Riesgos}.
			\item<3-> Plan de administración de \textbf{Costos}.
			\item<4-> Plan de administración de \textbf{Planificación}.
			\item<5-> Activos organizacionales del Proceso.
		\end{itemize}
	\end{column}
	\begin{column}{0.2\textwidth}
		\includegraphics[width=2cm]{img/input}
	\end{column}
\end{columns}
}

\frame
{
\frametitle{Análisis Cuantitativo de Riesgos}
\framesubtitle{Herramientas y Técnicas}
\begin{columns}
	\begin{column}{0.8\textwidth}
		\begin{itemize}
			\item<1-> \emph{Recolección} de datos y técnicas de representación.
			\item<2-> Análisis cuantitativo de riesgos y técnicas de modelamiento.
			\begin{itemize}
				\item<3-> Análisis de sensibilidad.
				\item<4-> Análisis de valores de \emph{expectación monetaria}.
				\item<5-> Modelamiento  y simulación.
			\end{itemize}
			\item<6-> Juicio Experto.
		\end{itemize}
	\end{column}
	\begin{column}{0.2\textwidth}
		\includegraphics[width=2cm]{img/tools}
	\end{column}
\end{columns}
}

\frame
{
\frametitle{Análisis Cuantitativo de Riesgos}
\framesubtitle{Salidas}
\begin{columns}
	\begin{column}{0.8\textwidth}
		Actualización de los registros de riesgos.
		\begin{itemize}
			\item<1-> Análisis \textbf{probabilístico} del proyecto.
			\item<2-> Probabilidad de \emph{obtener} los objetivos de costo y tiempo.
			\item<3-> Lista \textbf{priorizada} de riesgos cuantificados.
			\item<4-> \textbf{Tendencias} en los resultados cuantitativos de análisis de riesgos.
		\end{itemize}
	\end{column}
	\begin{column}{0.2\textwidth}
		\includegraphics[width=2cm]{img/output}
	\end{column}
\end{columns}
}


\section{Respuestas al plan de riesgos}
Las respuestas de planificación de riesgo deben ser adecuadas a la importancia del riesgo, el costo efectivo para enfrentarse al reto, realista en el contexto del proyecto, acordado por todas las partes involucradas, y es propiedad de una persona responsable. También debe ser oportuna, para seleccionar la mejor respuesta entre varias opciones de riesgo que se requieren a menudo.

Es el proceso de desarrollo de opciones y oportunidades para mejorar y reducir las amenazas a los objetivos del proyecto. Incluye la identificación y asignación de una persona (el dueño de respuesta al riesgo) para asumir la responsabilidad de cada uno de acuerdo a los riesgos y la respuesta de capitalización. Plan de respuestas al riesgo dirige los riesgos por su prioridad, inserta recursos y actividades en el presupuesto, controla cronograma y gestiona proyectos según sea necesario.

\subsection{Entradas}

\begin{itemize}
	\item Registro de Riesgos
	\item Plan de Administración de Riesgos
\end{itemize}

\subsection{Herramientas y Técnicas}

\begin{itemize}
	\item Estrategias para Riesgos Negativos o Amenazas
	\item Estrategias para Riesgos Positivos u Oportunidades
	\item Estrategias de Respuesta de Contingencia
	\item Juicio Experto
\end{itemize}

\subsection{Salidas}

\begin{itemize}
	\item Actualizaciones de Registro de Riesgos
	\item Decisiones de Contrato relacionadas a Riesgos
	\item Actualizaciones a Plan de Administración de Proyecto
	\item Actualizaciones al Documento de Proyecto
\end{itemize}


\section{Monitoreo y Control de Riesgos}
Monitorear y controlar los riesgos es el proceso de ejecución de planes de respuesta ante el riesgo, el seguimiento de los riesgos identificados, supervisar los riesgos residuales, identificando nuevos riesgos, y evaluar la eficacia del proceso de riesgo en todo el proyecto.

Respuestas de Planificación de riesgo que se incluyen en el plan de gestión de proyectos se ejecutan durante el ciclo de vida del proyecto, pero el trabajo del proyecto debe ser objeto de control permanente para las nuevos, cambiantes y actualizados riesgos.

El Monitor de Riesgos y Control de Procesos aplica técnicas, tales como la varianza y análisis de tendencias, que requieren el uso de información sobre el rendimiento generado durante la ejecución del proyecto:

\subsection{Entradas}

\begin{itemize}
	\item Registros de Riesgo
	\item Plan de Administración de Proyectos
	\item Información de Rendimiento en el Trabajo
	\item Reportes de Rendimiento
\end{itemize}

\subsection{Herramientas y Técnicas}

\begin{itemize}
	\item Revaloración de Riesgos
	\item Auditorías de Riesgos
	\item Varianza y Análisis de Tendencias
	\item Medición de Rendimiento Técnico
	\item Análisis de Reserva
	\item Reuniones de Estado
\end{itemize}

\subsection{Salidas}

\begin{itemize}
	\item Actualizaciones de Registro de Riesgos
	\item Actualizaciones de los Activos de Proceso de Organización
	\item Solicitudes de Cambio
	\item Actualizaciones al Plan de Administración de Proyecto
	\item Actualizaciones al Documento de Proyecto
\end{itemize}



\section{Conclusiones}
%APORTE RODRIGO:
La administración de los riesgos de un proyecto es un proceso que involucra a
todas las partes que o conforman, y que es muy necesario para poder estimar
los tiempos de desarrollo del mismo y de poder ayudar a evitar retrasos o el
fracaso del mismo. Además, las herramientas para poder realizar esto requieren de
personas que las manejen de buena forma y que puedan ver la visión holística
del problema para encontrar acciones de mitigación lo más adecuadas posibles.

%APORTE CRISTIAN:
Como equipo de trabajo tenemos muy clara, la importancia del todos los tópicos
que fueron repartidos entre los grupos de trabajos, pero creemos que la gestión
de los riesgos de un proyecto es fundamental, pues si no se tiene un especial cuidado
puedes llegar a acabar con un proyecto determinado.
En nuestra experiencia personal en la "Feria de Software", tuvimos la suerte de ser
los tres parte del mismo equipo, y al finalizar nos dimos cuenta que la parte más
dificil fué poder controlar los distintos riesgos de nuestro proyecto, tanto
del nivel técnico como a niver personal.

%APORTE GABRIEL:
Finalmente, se hace indispensable definir las respuestas que se pueden obtener al 
plantear el plan de riesgos, aplicando tanto herramientas, metodologías y 
técnicas, con el fin de manejar un conjunto de opciones a la hora de mitigar dichos
riesgos. Por otro lado todo proceso debe ser monitoreado y controlado, para ver 
que efectos tiene sobre el todo, con tal de hacer un seguimiento en el tiempo y
verificar que siga el camino establecido en los planes de gestión.

\end{document}
